\section{绪论}\label{sec:intro}

\subsection{背景：语言模型的上下文与状态}\label{sec:bg}

\textbf{先说循环网络这一路。}LSTM/GRU 及其变体\citep{hochreiter1997lstm,cho2014gru}
是最早把「固定大小的状态」写进架构的一路：
一步一递归，历史被压缩进一个形状固定的隐状态。这个想法本身是对的，但它没有走到今天，原因有两条。
其一，\textbf{状态太小}：一般只有 KB 到 MB 量级，容量撑不住长历史；其二，\textbf{效果不好}：
长程依赖学不动、训练也不稳定，于是被「更大的窗口 + 注意力」一路取代。
所以要把问题评准——症结不在「固定状态」这个想法，而在\textbf{当时的状态不够大、写入方式不够好}。
本文接下来做的事，就是把状态放大：我们研究的模型（Gridman/Ouro medium）仅 delta 规则矩阵记忆一项就有
$12{,}954{,}624$ 个 fp32 元素（约 $49.4$\,MiB），推理时 batch$=1$ 的四件套状态切片合计
$14{,}072{,}832$ 个 fp32 元素（约 $53.7$\,MiB：memory $49.4$ $+$ MTM $4.0$ $+$ queue $0.25$ $+$ recurrent $0.016$），
比典型 RNN 状态高出两个到四个数量级。

\textbf{既有架构只有一种信息载体：上下文窗口。}每一步预测都通过注意力读取窗口内所有 token；
窗口之外的历史模型完全看不见——除非把它重新写回窗口（重放），而重放的成本随长度增长
（注意力计算与显存都随长度线性膨胀）\citep{vaswani2017attention,dao2022flashattention,pope2023efficient}。
窗口越开越大是一条路\citep{beltagy2020longformer,ding2023longnet}，但它改变的是「一次能看多少」，
不是「能带多少」。

\textbf{持续状态是另一种载体。}一段形状固定的张量，跨上下文持续存在，在每次前向中被读写，
其容量\textbf{不随历史长度增长}：喂给它 1 千字节还是 100 万字节，状态大小不变。
窗口用「更大」换「记得更多」，状态用「固定」换「持续」——两者的张力正是本文的问题：
固定大小的状态到底能承担多少跨上下文的信息传递？这是一个可以用实验正面验证、给出定量结论的经验问题，
而不是「架构看起来很优雅」式的主张。

\textbf{全文围绕两点：}①\textbf{把状态做大}（scaling up）；
②\textbf{大状态带来的持续学习}——状态跨上下文持续被写入，模型在「读」的过程中仍在改变自己的行为，
而不是读完即定。逻辑链全文固定为：\emph{以前 RNN 效果不好 $\to$ 本文把状态放大 $\to$ 更大的状态带来持续学习。}

\subsection{问题：固定大小的持续状态能否承载跨上下文信息}\label{sec:problem}

\textbf{形式化。}一次前向写成
\begin{equation}
\mathrm{logits},\; s' = f(x, s; \mathit{lock}),
\label{eq:forward}
\end{equation}
其中 $x$ 是当前输入（本文实现中为一个 64 字节的 patch），$s$ 是状态——
\textbf{形状固定，与已处理的历史长度无关}，这就是「固定大小」的字面含义；$s'$ 是写入后的新状态；
$\mathit{lock}$ 是布尔开关，决定这一步 $s$ 是否接受新写入（$\mathit{lock}=\mathrm{true}$ 时状态只读，见 \S\ref{sec:model-io}）。
上下文的延续由 $s' = f(x,s)$ 在 patch 之间递归传递，不依赖 KV cache 的累积。

\textbf{两个待验证的子问题。}它们是本文的两条主张，第 \ref{sec:s1}、\ref{sec:s2} 章分别正面验证：

\begin{itemize}\itemsep4pt
\item \textbf{S1：状态是否持续携带跨上下文信息？}
\emph{成立条件}：\textbf{热启动显著优于空白}、\textbf{正常写入显著优于 $\mathit{lock}$}。
只有当前者成立，才说明状态里带着信息；只有当后者成立，才说明「往状态里写」这件事有因果贡献。
两条同时满足，S1 成立。
\begin{itemize}\itemsep2pt
\item \textbf{热启动（hot start）}：评测起点用训练收敛时保存的稳态状态快照（\S\ref{sec:train}），
      相当于「模型带着训练结束时的记忆来考试」。
\item \textbf{空白（blank）}：评测起点用全初始状态（恒等 memory + 零 recurrent/queue/MTM），
      相当于「模型失忆后开始读」。
\item \textbf{$\mathit{lock}$}：前向过程中状态只读不写——这段上下文的信息\textbf{不许}进入状态；
      它与「正常（commit）」构成一对受控对照，隔离出「写入」的因果贡献。
\end{itemize}
\item \textbf{S2：固定状态的有效上下文窗口有多大？}
\emph{成立条件}：\textbf{有效窗口显著大于 0}——固定状态在损失上确实等效于一段非零的滑动窗口。
\textbf{有效窗口}不是架构参数而是测出来的等效值\pending{测量方法见 \S\ref{sec:b3}，需滑窗基线}。
\end{itemize}

S1 问「有没有」（定性），S2 问「有多少」（定量）。第 \ref{sec:discussion} 章的汇总表把这两问与第 \ref{sec:s1}、\ref{sec:s2} 章的判定句直接对上。

\subsection{本文贡献}\label{sec:contrib}

\textbf{主线句}：以前 RNN 效果不好，是因为状态太小；本文把状态放大，大状态让「跨上下文携带信息」成立，
并进一步带来持续学习。三条贡献分别是这条链上的三段证据。

\textbf{贡献一：起点状态的系统消融结论（回答 S1，第 \ref{sec:s1} 章）。}\
五组对照——热启动 vs 空白（A1）、写入 vs $\mathit{lock}$（A2）、跨文档边界连续 vs 重置（A3）、
状态错配（A4）、组件级归因（A5/A6）——每组预先写明「状态在携带信息时结果应该长这样」的成立预期，
合起来逐条判定状态是否真在携带信息。这是「状态放大了，但里面到底有没有东西」的直接回答。
本文在该簇上的实测结论是：S1 \textbf{成立}，且携带者主要是块级循环状态（\S\ref{sec:s1-attr}）。

\textbf{贡献二：固定大小状态的有效窗口测量（回答 S2，第 \ref{sec:s2} 章）。}\
三把尺交叉定位：位置衰减（长文档上损失随位置怎么走，B1）、重复距离（自然重复内容隔多远还记着，B2）、
等效滑窗（把状态折算成一个窗口宽度数字，B3）。本文已完成前两把尺，第三把尺与精确检索预实验（B4）未完成，
如实标注。结论是\textbf{有效记忆距离约 $1$–$4$k token}：S2 的「显著大于 0」成立，
但远未达到「永续上下文」量级。

\textbf{贡献三：持续状态与有状态架构（第 \ref{sec:model} 章）。}\
架构层面的贡献有两层含义，这里给出的是\textbf{状态在前}的表述：
①\textbf{状态是微观实现、可以显式调整}——它可以整体读出、冻结（$\mathit{lock}$）、置换（重置/换入其他状态），
所以评测旋钮直接落在架构接口上，而不是像无状态模型那样只能从外部改提示词、改窗口；
②\textbf{记忆是宏观体现、可以多方面量度}——记忆不是外挂模块，就是四件套状态
（memory / recurrent / queue / MTM，见 \S\ref{sec:model-state}）在跨上下文尺度上表现出来的能力；
\textbf{持续学习是它在时间维度上的表现}——状态持续被写入，行为随之改变。

\draftnote{说明：原大纲把「teacher-forced bits/byte 评测协议」列为贡献之一，本稿按要求删除——
协议是支撑贡献一、二的手段，不单列为贡献；协议本身在第 \ref{sec:protocol} 章完整定义。}

\textbf{本文材料边界。}本文所报实验全部基于\textbf{单一检查点}（step 144000/705600，约 20\% 训练进度）的推理评测，
不做任何重训；已完成 A1–A3、A5、A6、B2 及其组件归因、外加稳定性（E 线）、
热启动记忆的领域分解与成员推断（D 线）、前向学习（G 线）；
未完成 A4、B3、B4、F 线（需重训）；受控遗忘曲线（C1）因实现缺陷作废，完整说明见附录 \ref{app:c1}。
